\section{Introduction}

Structure-based virtual screening has become a central tool in small-molecule
lead discovery, offering a route to interrogate chemical space at a scale that
is impractical with wet-lab assays alone. The explosion of make-on-demand
ultra-large libraries, which now span billions of purchasable molecules with
high reported synthesis success rates, has dramatically expanded the pool of
candidate ligands that can in principle be evaluated against a target
\citep{lyu2019ultralarge,tingle2023zinc22,irwin2005zinc}. Yet only a handful of
virtual screening campaigns deployed at this multi-billion scale have
translated into experimentally validated hits, in part because exhaustive
physics-based docking of entire libraries is computationally prohibitive
\citep{gorgulla2020virtualflow,gentile2022deepdocking}.

Several strategies have emerged to address this cost. Parallelized docking
platforms distribute calculations across high-performance computing clusters
\citep{gorgulla2020virtualflow,yaacoub2022ddgui}; active-learning and
Bayesian-optimization approaches train surrogate models to triage libraries so
that only the most promising fraction must be docked
\citep{graff2021accelerating,gentile2022deepdocking}; hierarchical pipelines
apply successively more expensive scoring stages to progressively smaller
compound pools; and GPU-accelerated docking programs further compress
wall-clock time. However, the ultimate impact of any of these strategies is
bounded by the accuracy of the underlying docking program, which must both
sample a biologically meaningful binding pose and score it in a way that
distinguishes true binders from decoys.

Leading proprietary docking programs such as Glide \citep{friesner2004glide}
and the grey-wolf-optimized variants integrated into VirtualFlow
\citep{gorgulla2020virtualflowgwo} offer strong benchmarked performance but
are not freely available for unrestricted academic use. AutoDock Vina
\citep{trott2010vina} is free and widely adopted but tends to trail
commercial programs on standard benchmarks. Recent deep-learning scoring
functions report competitive numbers on benchmarks such as CASF-2016
\citep{su2019casf2016,liu2017pdbbind,scantlebury2023pointvs}, yet their ability
to generalize to novel compounds and targets is limited, and data-leakage
between training and evaluation sets remains a concern even when stringent
similarity cutoffs are applied. When the binding site is known, as is typical
in virtual screening, physics-based methods continue to be the pragmatic
choice \citep{chaput2017usc,ebejer2019ligity,li2024cache}. Open-source physics
based platforms that combine accurate force fields, receptor flexibility, and
active-learning triage for billion-compound libraries remain scarce.

In this work, we introduce RosettaGenFF-VS and the RosettaVS protocol, built
on top of the Rosetta GALigandDock framework \citep{leman2020rosetta}. We
extend the underlying generalized force field with new atom types, new
torsional potentials, and an entropy model so that it can rank large pools of
compounds binding to the same target. We expose two docking modes: a fast
virtual screening express (VSX) mode for initial library triage and a virtual
screening high-precision (VSH) mode that restores full receptor sidechain
flexibility for final re-ranking of top candidates. These are wrapped by
OpenVS, an open-source AI-accelerated virtual screening platform that trains
a target-specific surrogate model on-the-fly to select promising compounds
for expensive docking, enabling routine screens of multi-billion compound
libraries on modest HPC resources. We benchmark the method on CASF-2016
\citep{su2019casf2016} and the Directory of Useful Decoys
\citep{mysinger2012dude}, then apply the platform to two unrelated
drug-discovery problems: a substrate receptor in the CUL2/CRL2 ubiquitin
ligase family, KLHDC2, which has recently attracted attention as a PROTAC
platform \citep{slusarz2023klhdc3}, and the voltage-sensing domain IV (VSD4)
of the human voltage-gated sodium channel NaV1.7, a validated but notoriously
difficult target for state-dependent small-molecule inhibitors
\citep{wang2015nav17methyl}. The KLHDC2 co-crystal structure determined for
the best initial hit validates the predicted binding pose, and the NaV1.7
leads display state-dependent inhibition with moderate selectivity over
NaV1.5 and hERG.

The contributions of this work are:
\begin{itemize}
  \item \textbf{Force field}: RosettaGenFF-VS extends Rosetta GALigandDock
  \citep{leman2020rosetta} with new atom types, torsional potentials, and an
  entropy term for ranking ligands binding to the same target.
  \item \textbf{Protocol}: RosettaVS supplies paired VSX (fast) and VSH
  (flexible-receptor, high-precision) docking modes for ultra-large screens.
  \item \textbf{Platform}: OpenVS is an open-source scalable virtual
  screening platform that combines these modes with active-learning triage,
  enabling multi-billion compound screens in under seven days on modest HPC
  resources.
  \item \textbf{Benchmarks}: We report leading CASF-2016 docking and
  screening power (top 1\% enrichment factor of 16.72 versus 11.9 for the
  second-best method) and DUD ROC enrichment, outperforming the second-best
  method by a factor of two at 0.5/1.0\% false-positive rates.
  \item \textbf{Applications}: Parallel KLHDC2 and NaV1.7 campaigns yielded
  seven low-micromolar KLHDC2 binders (14\% hit rate) and four
  single-digit-micromolar NaV1.7 antagonists (44.4\% hit rate), the best of
  which showed state-dependent hNaV1.7 inhibition with IC$_{50}$ = 1.32~\textmu M.
  \item \textbf{Structural validation}: A 2.0~\AA\ X-ray co-crystal structure
  of KLHDC2 with compound~29 validates the predicted binding pose of the
  computationally identified hit.
\end{itemize}

\section{Background and Related Work}

\paragraph{Ultra-large library docking.}
Make-on-demand libraries have expanded chemical space from the millions to
the billions of compounds \citep{lyu2019ultralarge,tingle2023zinc22}. Early
examples of ultra-large screens demonstrated that hit rates and chemotype
diversity can improve with increasing library size
\citep{lyu2019ultralarge,gorgulla2020virtualflow}, motivating the development
of platforms that can evaluate billions of molecules efficiently.
VirtualFlow demonstrates distributed ultra-large screening against targets
such as KEAP1 \citep{gorgulla2020virtualflow}, and later extensions
incorporate receptor flexibility via grey-wolf optimization
\citep{gorgulla2020virtualflowgwo}.

\paragraph{Active learning and deep-learning-guided screening.}
Rather than docking an entire library, active-learning approaches train a
cheap surrogate on a small fraction of docked compounds and use it to triage
the rest. Pool-based Bayesian optimization can recover a large share of the
top-scoring ligands in a 100-million-compound library after docking only a
few percent of it \citep{graff2021accelerating}. Deep Docking and its
successors deploy iterative neural-network surrogates to accelerate ligand
ranking in conjunction with a conventional docker
\citep{gentile2022deepdocking,yaacoub2022ddgui}. Recent community exercises
such as CACHE suggest that while such acceleration is reliable, hit rates on
truly novel pockets remain modest, underscoring the importance of the
underlying scoring function \citep{li2024cache}.

\paragraph{Docking programs and scoring functions.}
Glide uses hierarchical conformational search with an empirical-plus-force-
field scoring function and reports strong redocking accuracy on standard
benchmarks \citep{friesner2004glide}. AutoDock Vina introduced a smoothed
empirical scoring function with efficient optimization and multithreading,
and remains a widely used open-source baseline \citep{trott2010vina}. The
PDBbind and CASF benchmarks quantify scoring, ranking, docking, and
screening power of scoring functions on diverse protein-ligand complexes
\citep{liu2017pdbbind,su2019casf2016}, while the Directory of Useful Decoys
and its enhanced version evaluate virtual-screening discrimination using
property-matched decoys \citep{mysinger2012dude}. Consensus docking analyses
and knowledge-based ligand methods illustrate both the strengths and the
failure modes of current approaches \citep{chaput2017usc,ebejer2019ligity}.

\paragraph{Machine-learning scoring functions.}
Neural-network scoring functions have shown competitive CASF numbers, yet
concerns about train/test leakage persist even when ligand Tanimoto and
protein sequence-identity cutoffs are applied. Recent work argues that
careful dataset construction and interpretability testing are needed before
these methods can be trusted as drop-in replacements for physics-based
scoring \citep{scantlebury2023pointvs}. Additionally, docking apo protein
structures typically degrades virtual-screening performance, motivating
flexible-receptor refinement and holo-like pocket conformation generation
\citep{guterres2020apo}.

\paragraph{KLHDC2 and NaV1.7 as test targets.}
KLHDC2 is a substrate receptor of the CUL2/CRL2 E3 ubiquitin ligase complex
and recognizes C-terminal diglycine degrons with high affinity
\citep{slusarz2023klhdc3}. Its tractable, pocket-like degron-binding site
has made it a candidate platform for PROTAC-style targeted protein
degradation. The human voltage-gated sodium channel NaV1.7, in turn, is a
well-validated analgesic target for which subtype-selective small-molecule
inhibitors have been pursued through interactions with the fourth
voltage-sensing domain (VSD4); subtle chemical modifications of known aryl
sulfonamide scaffolds can introduce new modes of gating modulation, making
the domain a rewarding but structurally subtle pocket for virtual screening
\citep{wang2015nav17methyl}. We chose these two targets to test the
generality of RosettaVS across protein folds, pocket physicochemistry, and
assay modalities.

\paragraph{This work.}
Against this backdrop, RosettaGenFF-VS, RosettaVS, and OpenVS combine an
open-source, physics-based scoring function with flexible-receptor docking
and active-learning triage. Unlike deep-learning scoring functions, the
method does not rely on the distribution of training ligands; unlike
proprietary commercial pipelines, all components are free and open-source;
and unlike previous active-learning approaches, receptor flexibility is
modeled at the final VSH stage. The benchmark and experimental results that
follow justify each of these choices.
